\subsection{Interval Notation: Unions of Intervals} 
Some sets are made up of \makeblank interval.
\begin{Example} The set shown below is made up of three intervals.
\begin{icpic}
\numlinec{7}
\graphoc{-5}{-2}
\graphoo{1}{3}
\graphca{4}{8}
\end{icpic}
The intervals are:
\begin{enumerate}
\item \makeblank
\item \makeblank
\item \makeblank
\end{enumerate}
The entire set is the \makeblank[1in] of the intervals, (Note that the itervals do not {\makeblank}!) To write this set in interval notation, we write:
\[
\makeblank[3in]
\]
\end{Example}
\begin{Note} We only use this method when our set has \makeblank[1in] in it, so there are two or more pieces that don't \makeblank.
\end{Note}
\begin{Example}
Rewrite $\bigl(2,7\bigr)\cup\bigl[5,8\bigr]$ as a single interval.
\begin{cpic}{}{3}
\end{cpic}
\end{Example}